\subsection{Absolute values on Q}

PROPOSITION 4. Let \(f\) be a mapping of \(\mathbf{Q}\) to \(\mathbf{R}_+\) belonging to \(\mathcal{V}(\mathbf{Q})\) . Then:

\begin{enumerate}
\def\labelenumi{(\roman{enumi})}
\item
  Either \(f\) is the improper absolute value on \(\mathbf{Q}\) .
\item
  Or there exists a real number \(a\) and a prime number \(p\) such that \(0 < a < 1\) and \(f = a^{v_p}\) , where \(v_p\) is the \(p\) -adic valuation.
\item
  Or there exists \(s > 0\) such that \(f(x) = |x|^s\) for all \(x \in \mathbf{Q}\) .
\end{enumerate}

In case (iii) for \(f\) to be an absolute value on \(\mathbf{Q}\) , it is necessary and sufficient that \(0 < s \leqslant 1\) .

Suppose first that \(f(n) \leqslant 1\) for every integer \(n > 0\) . By Proposition 3 of no. 2 there exist a real number \(b\) and a valuation \(v\) on \(\mathbf{Q}\) such that \(0 < b < 1\) and \(f = b^v\) . Now, we know (§ 3, no. 4, Example 4) that the only valuations on \(\mathbf{Q}\) are (up to equivalence) the improper valuation and the \(p\) -adic valuations \(v_p\) ; we therefore have either case (i) or case (ii).

Suppose from now on that there exists an integer \(h > 0\) such that \(f(h) > 1\) ; by no. 1, Corollary 2 to Proposition 2, there exists a number \(\rho > 0\) such that \(f^{\rho}\) is an absolute value; let us write

\[
g (x) = \rho \log (f (x)) / \log | x |
\]

for every rational number \(x \neq 0\) . Let \(a, b\) be two integers \(\geqslant 2\) ; for every integer \(n \geqslant 2\) let \(q(n)\) denote the integral part of \(n\) . log \(a / \log b\) , in other words the least integer \(m\) such that \(a^n < b^{m+1}\) ; the expansion of \(a^n\) to base \(b\) is therefore

\[
a ^ {n} = c _ {0} + c _ {1} b + \dots + c _ {q (n)} b ^ {q (n)}\tag{2}
\]

where \(0 \leqslant c_{i} < b\) for \(0 \leqslant i \leqslant q(n)\) . As \(f^{\rho}\) is an absolute value, \(f^{\rho}(c_{i}) \leqslant c_{i} \leqslant b\) and we therefore deduce from (2) that

\[
\begin{array}{r l} (f (a)) ^ {n \rho} = (f (a ^ {n})) ^ {\rho} & \leqslant b (1 + (f (b)) ^ {\rho} + \dots + (f (b)) ^ {\rho q (n)}) \\ & \leqslant b (q (n) + 1) (\sup (1, (f (b)) ^ {\rho})) ^ {q (n)}. \end{array}
\]

Taking logarithms on both sides of this inequality and dividing by \(n.\log a\) , we obtain

\[
g (a) \leqslant \frac {\log b}{n . \log a} + \frac {\log (q (n) + 1)}{q (n)} \cdot \frac {q (n)}{n . \log a} + \frac {\sup (0 , \rho \log f (b))}{\log a} \cdot \frac {q (n)}{n}.\tag{3}
\]

Note now that as \(n\) tends to \(+\infty\) , \(q(n) / n\) tends to \(\log a / \log b\) ; therefore \(q(n)\) tends to \(+\infty\) and

\[
\log (q (n) + 1) / q (n)
\]

tends to 0 (Functions of a Real Variable, Chapter III, § 2, no. 1). Taking the limit in (3), we obtain

\[
g (a) \leqslant \frac {\sup (0 , \rho \log f (b))}{\log b} = \sup (0, g (b)).\tag{4}
\]

But \(f(h) > 1\) , whence \(g(h) > 0\) ; if \(a\) is replaced by \(h\) in (4), we obtain \(\sup(0, g(b)) > 0\) and hence

\[
\sup (0, g (b)) = g (b).
\]

Then, for any integers \(a\) , \(b\) at least equal to 2, \(g(a) \leqslant g(b)\) and therefore \(g(a) = g(b)\) , exchanging the roles of \(a\) and \(b\) . In other words, there exists a constant \(\lambda\) such that \(g(a) = \lambda\) for every integer \(a \geqslant 2\) ; if we write \(s = \lambda / \rho\) , then \(f(a) = |a|^s\) for every integer \(a \geqslant 2\) . As \(f(xy) = f(x)f(y)\) and \(f(-x) = f(x) = |x|^s\) for all \(x \in \mathbf{Q}\) . Finally, if \(0 < s \leqslant 1\) , we know that \(x \mapsto |x|^s\) is an absolute value (General Topology, Chapter IX, §3, no. 2); conversely, if \(s\) is such that \(x \mapsto |x|^s\) is an absolute value on \(\mathbf{Q}\) , then \((1 + 1)^s \leqslant 1^s + 1^s\) , that is \(2^s \leqslant 2\) , whence \(s \leqslant 1\) .

